% 图契约：本题的源、边界和推导逐项对应；箭头符号按图注解释。
\documentclass[tikz,border=7pt]{standalone}
\usepackage{fontspec,xeCJK,amsmath,amssymb}
\IfFontExistsTF{Times New Roman}{\setmainfont{Times New Roman}}{\setmainfont{TeX Gyre Termes}}
\IfFontExistsTF{Heiti SC}{\setCJKmainfont[AutoFakeBold=1.5]{Heiti SC}}{\setCJKmainfont[AutoFakeBold=1.5]{FandolHei-Regular}}
\usetikzlibrary{arrows.meta,bending,calc,patterns,positioning}
\definecolor{ink}{HTML}{243745}
\definecolor{blue}{HTML}{275DAB}
\definecolor{green}{HTML}{177C63}
\definecolor{red}{HTML}{BD4935}
\definecolor{orange}{HTML}{AA6C16}
\definecolor{purple}{HTML}{715393}
\tikzset{every node/.style={font=\fontsize{9.3}{12.5}\selectfont,text=ink},
axis/.style={-{Stealth[length=4pt]},line width=.65pt,draw=ink!65},
vec/.style={-{Stealth[length=5pt]},line width=1pt,draw=blue},
short/.style={-{Stealth[length=3pt]},line width=.8pt},
guide/.style={dashed,draw=ink!40,line width=.5pt},
card/.style={rounded corners=2pt,draw=blue!65,fill=blue!4,align=center,text width=6.5cm,inner xsep=9pt,inner ysep=7pt},
title/.style={font=\fontsize{10.5}{14}\selectfont\bfseries,text width=15cm,align=center},
label/.style={font=\fontsize{8.4}{11}\selectfont,align=center},
note/.style={label,text width=6.2cm}}
\begin{document}
\begin{tikzpicture}[x=1cm,y=1cm]
\path[use as bounding box] (-8,-3.5) rectangle(8,3.5);
\node[title] at (0,3.15) {M3-010\quad 介质球壳中的三种高斯定律};
\begin{scope}[shift={(-4,0)}]
\filldraw[draw=ink!55,fill=blue!7] (0,0)circle(1.8);
\filldraw[draw=ink!55,fill=white] (0,0)circle(.9);

\draw[guide] (0,0)--(110:.9) node[pos=.8,left] {$a$};
\draw[guide] (0,0)--(40:1.8) node[pos=.8,above left] {$b$};
\node[label] at(0,1.4){$\varepsilon_r$};\node[label] at(-1.9,2){真空};
\draw[vec,green] (0:.95)--(0:1.65) node[midway,above] {$\mathbf P$};
\draw[vec,orange] (180:.9)--(180:.4);
\node[label,orange] at(-1.05,.28){$\mathbf n_a$};
\draw[vec,orange] (20:1.8)--(20:2.35) node[above] {$\mathbf n_b$};
\node[label,blue] at(-1.65,-.6){$\sigma_{b,a}<0$};
\node[label,red] at(1.4,-1.9){$\sigma_{b,b}>0$};
\fill[red] (0,0) circle(2.3pt) node[below=3pt] {$Q$};
\draw[green,dashed,thick] (0,0)circle(1.25);\node[label,green,fill=white] at(-1.25,1.05){$S_r$};\node[label] at(0,-2.55){符号示意取 $Q>0,\ \varepsilon_r>1$};\end{scope}
\node[card,draw=blue!65,fill=blue!4] at (3.9,1.8) {$\Phi_D=Q$：三个区域均相同};
\node[card,draw=green!65,fill=green!4] at (3.9,0) {$a<r<b$：$q_b=-Q\alpha$\\$\varepsilon_0\Phi_E=Q/\varepsilon_r,\ \Phi_P=Q\alpha$};
\node[card,draw=orange!65,fill=orange!4] at (3.9,-1.8) {$r<a$ 或 $r>b$：$q_b=0$\\$\varepsilon_0\Phi_E=Q,\ \Phi_P=0$};
\end{tikzpicture}\end{document}